% Part: sets-functions-relations
% Chapter: arithmetization
% Section: cauchy

\documentclass[../../../include/open-logic-section]{subfiles}

\begin{document}

\olfileid{sfr}{arith}{cauchy}

\olsection{ضمیمہ: حقیقی عدد بطور کوشی لڑیاں}

\olref[cuts]{sec} وچ اسی حقیقی عدداں نوں ڈیڈیکائنڈ کٹاں دے طور تے
بنایا سی۔ ایس حصے وچ اسی اک متبادل بناوٹ سمجھاؤندے آں۔ ایہ کوشی دی
(اوس شے دی جیہنوں ہُن اسی آکھدے آں) کوشی لڑی والی تعریف اُتے ٹکی اے؛
پر ایس تعریف نال حقیقی عدد \emph{بنان} دا کم انیسویں صدی دے ہور مصنفا،
خاص طور تے وائرشٹراس، ہائن، مے رے تے کانتور، دا اے۔ (چنگی تاریخی روداد
لئی \citeauthor{OConnorRobertson:RN} \citeyear{OConnorRobertson:RN} ویکھو۔)

انیسویں صدی ول جان توں پہلاں سائمن سٹیون (1548--1620) اُتے غور کرنا
فائدے مند اے۔ مختصر ایہ کہ سٹیون نے سمجھیا کہ اسی ہر حقیقی عدد نوں
اوہدی اعشاری پھیلاؤ راہیں سوچ سکدے آں۔ ایس لئی حتیٰ کہ $\sqrt{2}$
ورگے غیر ناطق عدد دا وی اک سوہنا اعشاری پھیلاؤ اے، جیہڑا ایویں شروع
ہوندا اے:
\[
	1.41421356237\ldots
\]
سیٹ تھیوری وچ اعشاری پھیلاؤ دا نمونہ بنانا بہت سوکھا اے: اوہناں نوں
بس فنکشن $d \colon \Nat \to \Nat$ سمجھو، جتھے $d(n)$ اوہ $n$واں
اعشاری مقام اے جیہدے وچ ساڈی دلچسپی اے۔ فیر اعشاری نقطے توں پہلاں
آؤندے حقیقی عدد دے حصے (ایتھے صرف $1$) نوں سنبھالن لئی اک چھوٹی تبدیلی
لوڑیں دی اے۔ مثال لئی $0.999\ldots = 1$ دی ضمانت لئی اک ہور تبدیلی
(اک ہم ارزیتی تعلق) وی لوڑیں دی اے۔ پر سٹیون دے ڈھنگ نال سیٹ تھیوری
اندر حقیقی عدداں دی پوری طرح سخت بناوٹ دینا اوکھا نہیں۔

سٹیون ساڈا مرکزی موضوع نہیں۔ (سٹیون بارے ہور لئی
\citealt{KatzKatz2012} ویکھو۔) پر ایتھے اک بہت نیڑا خیال اے۔
$\sqrt{2}$ دے اعشاری پھیلاؤ نوں سدھا لین دی بجائے اسی ودھدی درستی
والے پھیلاؤ لے کے، $\sqrt{2}$ دیاں ہور تے ہور درست ناطق تقریبیاں دی
اک \emph{لڑی} اُتے غور کر سکدے آں:
\[
1, 1.4, 1.414, 1.4142, 1.41421,\ldots
\]
حقیقی عدداں نوں ”ہور تے ہور چنگیاں تقریبیاں“ راہیں سمجھ سکّن دا خیال
سانوں سوجھاں دی اک ہور لڑی دا مُڈھ دیندا اے (اوہناں سوجھاں ورگی
جیہڑیاں اسی $\Rat$ دی $\Int$ توں، یا $\Int$ دی $\Nat$ توں بناوٹ
ویلے ورتیاں سن)۔ بنیادی بصیرتاں ایہ نیں:
\begin{enumerate}
	\item ہر حقیقی عدد نوں اک (شاید لامتناہی) اعشاری پھیلاؤ دے طور تے لکھیا جا سکدا اے۔
	\item اک (شاید لامتناہی) اعشاری پھیلاؤ وچ بند معلومات نوں ناطق عدداں دی لڑی وی اوہناں ہی پورا بند کر سکدی اے۔
	\item ناطق عدداں دی لڑی نوں $\Nat$ توں $\Rat$ ول فنکشن سمجھیا جا سکدا اے؛ بس $f(n)$ نوں لڑی دا $n$واں ناطق عدد لو۔
\end{enumerate}
بے شک $\Nat$ توں $\Rat$ ول ہر \emph{کوئی وی} فنکشن سانوں حقیقی
عدد نہیں دیوے گا۔ مثال لئی ایہ فنکشن ویکھو:
\[
	f(n) =\begin{cases}
			1 & \text{جے }n\text{ طاق اے}\\
			0 &\text{جے }n\text{ جفت اے}
		\end{cases}
\]
اصل فکر ایہ اے کہ لڑی $0,1,0,1,0,1,0,\ldots$ کسے حقیقی عدد ول
”سمٹدی“ نہیں لگدی۔ سو اوہ لڑیواں لین دی ضمانت لئی جیہڑیاں کسے حقیقی
عدد ول سمٹدیاں نیں، سانوں اپنی توجہ اوہناں لڑیواں تک محدود کرنی پینی
اے جیہناں دی کوئی \emph{حد} ہووے۔

اسی \olref[his][set][limits]{sec} وچ حد دا خیال پہلے ای ویکھ چکے آں۔
پر اسی عین اوہی تعریف نہیں ورت سکدے جیہڑی اوتھے ورتی سی۔ اوتھے اظہار
”$(\forall \epsilon>0)$“ وچ لُکّویں طور تے حقیقی عدداں اُتے مقداری
بندھن سی؛ نالے اسی حقیقی عدداں اُتے فنکشناں دیاں حداں اُتے غور کر رہے
ساں۔ سو اوس تعریف نوں ورتنا حقیقی عدداں نوں پہلے توں ہی من لینا ہوندا؛
تے اسی عین اوہناں نوں \emph{بنان} لگے آں۔ خوش قسمتی نال اسی حد دا اک
بہت نیڑا خیال ورت سکدے آں۔
\begin{defn}\ollabel{def:CauchySequence}
  فنکشن $f: \Nat \to \Rat$ اک \emph{کوشی لڑی} اُدوں تے صرف اُدوں
  اے جدوں ہر مثبت $\epsilon \in \Rat$ لئی
  $(\exists \ell \in \Nat)(\forall m, n > \ell)|f(m) - f(n)| < \epsilon$ ہووے۔
\end{defn}

حد دا عام خیال پہلے والا ای اے: جے تسیں درستی دی اک خاص سطح
(جیہڑی~$\epsilon$ نال نپی اے) چاہندے او، تاں ویکھن لئی اک ”علاقہ“
(~$\ell$ توں وڈا ہر انپٹ) اے۔ نالے ایہ ویکھنا سوکھا اے کہ ساڈی لڑی
$1$، $1.4$، $1.414$، $1.4142$، $1.41421$\ldots دی اک حد اے: جے تسیں
$\sqrt{2}$ دی تقریب $\nicefrac{1}{10^{n}}$ توں گھٹ غلطی نال چاہندے او،
تاں بس $n$ویں توں بعد دا کوئی وی رکن ویکھ لو۔

فیر صاف خیال ایہ ہووے گا کہ حقیقی عدد بس کوئی وی کوشی لڑی
\emph{اے}۔ پر $\Int$ تے $\Rat$ دیاں بناوٹاں وانگ ایہ گل بہت سادی
ہووے گی: ہر مقرر حقیقی عدد نوں کئی وکھ وکھ کوشی لڑیاں دَسدیاں نیں۔
ایہ ویکھن دا اک سادہ طریقہ ایہ اے۔ کوئی کوشی لڑی~$f$ دِتی ہووے تاں
$g$ نوں عین $f$ والا فنکشن تعریف کرو، بس $g(0)\neq f(0)$ ہووے۔ کیوں
جو دونیں لڑیاں پہلے عدد توں بعد ہر تھاں رلدیاں نیں، اسی آخرکار ایہ
آکھنا چاہواں گے کہ \olref{def:CauchySequence} وچ ورتے معنی مطابق
اوہناں دی حد اکّو اے، تے ایس لئی اوہناں نوں اکّو حقیقی عدد ”تعریف
کردیاں“ سمجھنا چاہیدا اے۔ سو سانوں واقعی ایہناں کوشی لڑیواں نوں اکّو
حقیقی عدد سمجھنا چاہیدا اے۔

ایس لئی سانوں فیر کوشی لڑیواں اُتے اک ہم ارزیتی تعلق تعریف کرنا، تے
حقیقی عدداں دی ہم ارزیت \emph{طبقیاں} نال شناخت کرنی پینی اے۔ پہلے
سانوں اوس فنکشن دا خیال چاہیدا اے جیہڑا حد وچ $0$ ول جاندا اے۔ ہر
فنکشن $h : \Nat \to \Rat$ لئی آکھو کہ \emph{$h$ حد وچ $0$ ول
جاندا اے} اُدوں تے صرف اُدوں جدوں ہر مثبت $\epsilon \in \Rat$ لئی
$(\exists \ell \in \Nat)(\forall n > \ell)|h(n)| <
\epsilon$ ہووے۔\footnote{ایہدی \olref[his][set][limits]{sec} وچ
$\lim_{x \mathord{\rightarrow}\infty}f(x) = 0$ دی تعریف نال
تُلا کرو۔} ہور، جتھے $f$ تے $g$، $\Nat \to \Rat$ فنکشن نیں،
مقرر کرو کہ $(f-g)(n) = f(n) - g(n)$۔ ہُن تعریف کرو:
\[
	f \Realequiv g \text{ اُدوں تے صرف اُدوں جدوں $(f-g)$ حد وچ $0$ ول جاندا اے}
.
\]
سانوں جانچنا پینا اے کہ $\Realequiv$ اک ہم ارزیتی تعلق اے؛ تے ایہ
اے۔ فیر جے اسی چاہواں، تاں حقیقی عدداں نوں $\Realequiv$ تحت، $\Nat
\to \Rat$ دیاں ساریاں کوشی لڑیواں دے ہم ارزیت طبقے تعریف کر سکدے آں۔
% PNBARI-009 ماخذ دی درستی: حقیقی عدداں دی شناخت ہم ارزیتی تعلق نال نہیں، اوس تعلق تحت کوشی لڑیواں دے طبقیاں نال ہوندی اے۔

\begin{prob}
$f(n) = 0$، ہر $n$ لئی، لو۔ $g(n) = \frac{1}{(n+1)^2}$ لو۔ وکھاؤ
کہ دونیں کوشی لڑیاں نیں، تے واقعی دونیں فنکشناں دی حد $0$ اے، ایس لئی
$f \sim_\Real g$ وی اے۔
\end{prob}

ایہ کر لین توں بعد اسی عام وانگ اوس ہم ارزیت طبقے لئی
$\equivrep{f}{\Realequiv}$ لکھاں گے جیہدے وچ $f$، !!a{element} اے۔
پر لکھت نوں پڑھّن جوگ رکھن لئی اگے اسی زیریں اشاریہ چھڈ کے صرف
$\equivrep{f}{}$ لکھاں گے۔ نالے اسی مقرر کردے آں کہ ہر $q \in \Rat$
لئی $q_{\Real} = \equivrep{c_{q}}{}$، جتھے $c_{q}$ مستقل فنکشن اے،
یعنی $c_q(n) = q$، ہر $n \in \Nat$ لئی۔ فیر اسی حقیقی عدداں اُتے
بنیادی تعلق تے عمل تعریف کردے آں، مثال لئی:
\begin{align*}
	\equivrep{f}{} + \equivrep{g}{} &= 	\equivrep{(f + g)}{} \\
	\equivrep{f}{} \times 	\equivrep{g}{} &= \equivrep{(f \times g)}{}
\end{align*}
جتھے $(f + g)(n) = f(n) + g(n)$ تے $(f \times g)(n) = f(n) \times
g(n)$۔ بے شک سانوں ایہ وی جانچنا پینا اے کہ $(f + g)$، $(f-g)$ تے
$(f\times g)$ وچوں ہر اک وی کوشی لڑی اے، جدوں $f$ تے $g$ کوشی لڑیاں
ہون؛ پر اوہ نیں، تے اسی ایہ کم تہاڈے لئی چھڈدے آں۔

آخر وچ اسی ترتیب دا اک تصور تعریف کردے آں۔ آکھو کہ
$\equivrep{f}{}$ \emph{مثبت} اے اُدوں تے صرف اُدوں جدوں
$\equivrep{f}{}\neq 0_\Real$ تے
$(\exists \ell \in \Nat)(\forall n > \ell)0 < f(n)$ دونیں ہون۔
فیر آکھو کہ $\equivrep{f}{} < \equivrep{g}{}$ اُدوں تے صرف اُدوں
جدوں $\equivrep{(g - f)}{}$ مثبت اے۔ سانوں جانچنا پینا اے کہ ایہ
صحیح طور تے تعریف شدہ اے (یعنی ایہ ہم ارزیت طبقے وچوں ”نمائندہ“ فنکشن
دی چُون اُتے منحصر نہیں)۔
% PNBARI-010 ماخذ دی درستی: حقیقی ہم ارزیت طبقے دی غیر صفری شرط اندر صحیح ٹائپ والا صفر 0_Real اے، 0_Rat نہیں۔
%\begin{align*}
%	\equivrep{f}{} \leq \equivrep{g}{} \text{ اُدوں تے صرف اُدوں جدوں }&\equivrep{f}{} = \equivrep{g}{} \text{ یا }\\
%	&(\exists \ell \in \Nat)(\forall n \in \Nat)(\ell < n \lif f(n) \leq g(n))
%\end{align*}
پر ایہ کر لین توں بعد ایہ وکھاؤنا کافی سوکھا اے کہ ایہ تعریفاں درست
الجبری خاصیتاں دیندیاں نیں؛ یعنی:
\begin{thm}\ollabel{thm:cauchyorderedfield}
کوشی لڑیواں دے ہم ارزیت طبقے اک ترتیب دار میدان بناؤندے نیں۔
% PNBARI-011 ماخذ دی درستی: ایتھے، اگلی مشق تے تکمیلیت دے قضیے وچ میدان دیاں ہستیاں خام کوشی لڑیاں نہیں بلکہ اوہناں دے ہم ارزیت طبقے نیں؛ اگلے ثبوت وچ نمائندہ لڑیاں والا نشان محفوظ اے۔
\end{thm}

\begin{proof}
مشق۔
\end{proof}

\begin{prob}
ثابت کرو کہ کوشی لڑیواں دے ہم ارزیت طبقے اک ترتیب دار میدان بناؤندے نیں۔
\end{prob}

ایویں بنائے حقیقی عدداں کول تکمیلیت دی خاصیت اے، ایہ ثابت کرنا ہور
اوکھا اے، سو اسی ثبوت دیندے آں۔

\begin{thm}
کوشی لڑیواں دے ہم ارزیت طبقیاں دا ہر غیر خالی سیٹ، جیہدی بالائی حد
ہووے، اک سب توں چھوٹی بالائی حد رکھدا اے۔
\end{thm}

\begin{proof}[ثبوت دا خاکہ] $S$ اوہناں کوشی نمائندہ لڑیواں دا کوئی وی
غیر خالی سیٹ ہووے جیہدے ہم ارزیت طبقیاں دی اک بالائی حد اے۔ سو کوئی
$p \in \Rat$ اے جیہدے لئی $p_{\Real}$، $S$ دے نمائندہ طبقیاں دی
بالائی حد اے۔ $r \in S$ لو؛ فیر کوئی $q \in \Rat$ اے جیہدے لئی
$q_{\Real} < r$۔ سو جے $S$ دی سب توں چھوٹی بالائی حد موجود اے، تاں
اوہ $q_\Real$ تے $p_\Real$ دے وچکار (دونیں سرے شامل کر کے) اے۔

اسی بازگشتی نصف بندی نال سب توں چھوٹی بالائی حد ول اکّو ویلے ہِٹھاں
تے اُتّوں نزدیک آ کے سمٹاں گے۔ خاص طور تے اسی دو فنکشن $f, g \colon \Nat \to \Rat$
تعریف کردے آں؛ مقصد ایہ اے کہ $f$ اُتّوں اوس حد ول سمٹے، تے $g$ ہِٹھاں
توں اوس ول سمٹے۔ اسی ایہ تعریفاں توں شروع کردے آں:
\begin{align*}
	f(0) &= p \\
	g(0) &= q
\end{align*}
فیر جتھے $a_n = \frac{f(n) + g(n)}{2}$، مقرر کرو:\footnote{ایہ اک
بازگشتی تعریف اے۔ پر اسی ہُن \emph{تک} ایہ سوچن دی کوئی وجہ نہیں
دِتی کہ بازگشتی تعریفاں درست نیں۔}
\begin{align*}
	f(n+1) &=
	\begin{cases}
		a_n &\text{جے }(\forall h \in S)\equivrep{h}{} \leq (a_n)_\Real\\
		f(n)&\text{نہیں تاں}
	\end{cases}\\
	g(n+1) &=
	\begin{cases}
		a_n &\text{جے }(\exists h \in S)\equivrep{h}{} \geq (a_n)_\Real\\
	 	g(n) &\text{نہیں تاں}
	\end{cases}
\end{align*}
$f$ تے $g$ دونیں کوشی لڑیاں نیں۔ (ایہ کافی سوکھ نال جانچیا جا سکدا
اے، پر اسی ایہنوں مشق لئی چھڈدے آں۔) نوٹ کرو کہ فنکشن $(f-g)$ حد وچ
$0$ ول جاندا اے، کیوں جو ہر قدم اُتے $f$ تے $g$ وچلا فرق آدھا ہو جاندا
اے۔ ایس لئی $\equivrep{f}{} = \equivrep{g}{}$۔

پہلاں اسی وکھاؤندے آں کہ $\equivrep{f}{}$، $S$ دی بالائی حد اے، یعنی
$(\forall h \in S)\equivrep{h}{} \leq \equivrep{f}{}$۔
(اگے ودھدے اسی \olref{thm:cauchyorderedfield} نوں ورتاں گے۔) $h \in S$
لو تے تضاد لئی فرض کرو کہ $\equivrep{f}{} < \equivrep{h}{}$، سو
$0_\Real < \equivrep{(h-f)}{}$۔ کیوں جو $f$ اک یک رخا گھٹدی کوشی
لڑی اے، کوئی $n \in \Nat$ اے جیہدے لئی
$\equivrep{(c_{f(n)} - f)}{} < \equivrep{(h-f)}{}$۔ سو:
\[
	(f(n))_\Real = \equivrep{c_{f(n)}}{} < \equivrep{f}{} + \equivrep{(h-f)}{} = \equivrep{h}{},
\]
پر ایہ اوس حقیقت دے خلاف اے کہ بناوٹ مطابق
$\equivrep{h}{} \leq (f(n))_\Real$۔

اگے اسی وکھاؤندے آں کہ $\equivrep{f}{} = \equivrep{g}{}$، $S$ دی
\emph{سب توں چھوٹی} بالائی حد اے۔ سو $j$ کوئی وی کوشی لڑی ہووے تے
فرض کرو کہ $\equivrep{j}{} < \equivrep{g}{}$۔ اوپر وانگ دلیل نال
(ایہ حقیقت ورت کے کہ $g$ \emph{ودھدی} اے)، کوئی $n \in \Nat$ اے
جیہدے لئی $\equivrep{j}{} < (g(n))_\Real$۔ پر بناوٹ مطابق کوئی
$h \in S$ اے جیہدے لئی $(g(n))_\Real \leq \equivrep{h}{}$، سو
$\equivrep{j}{} < \equivrep{h}{}$، تے ایس لئی $\equivrep{j}{}$،
$S$ دی بالائی حد نہیں۔
\end{proof}

\end{document}
